% SECCIÓN 3: LA LEX ARTIS MÉDICA — NATURALEZA, FUNCIÓN Y CRITERIOS DE DETERMINACIÓN
%
% Objetivos de esta sección:
%   1. Definir la lex artis médica como el estándar de conducta exigible al profesional
%      de salud en el ejercicio de su actividad.
%   2. Distinguir:
%      - Lex artis general: estándar de la especialidad médica (objetivado).
%      - Lex artis ad hoc: estándar ajustado a las circunstancias del caso concreto
%        (paciente, recursos disponibles, urgencia).
%   3. Analizar cómo el incumplimiento de la lex artis configura la antijuridicidad
%      y el factor de atribución (culpa profesional).
%   4. Revisar los criterios usados en la jurisprudencia comparada (España, Argentina)
%      para la determinación de la lex artis.
%   5. Identificar problemas probatorios: prueba pericial, historia clínica, protocolos.
%
% Referencias clave: Fernández Cruz, Galán Cortés, Lorenzetti, CAS. relevantes
% Extensión estimada: 1.500–2.000 palabras

\section{La \textit{lex artis} médica: naturaleza, función y criterios de determinación}
\label{sec:lex_artis}

